\chapter{HotChocolate 14, in Full}
\label{app:hc14}

\section{Scope and Migration Map}

This is the single-service compatibility example used by the
\enquote{On HotChocolate 14} callouts. It keeps sessions and speakers in one
database. It is not a port of the final four-service graph, and missing
services here are not evidence that the older library cannot implement them.
Use a separate example directory; its Sessions service uses port 5001.

The older Hot Chocolate release no longer receives security fixes.
Appendix~\ref{app:toolchain} records the tested versions and the current
security-policy source. For an existing application, these differences are
migration aids, not reasons to remain on that release.

\begin{center}
\small
\begin{tabular}{p{0.40\linewidth}p{0.47\linewidth}}
\toprule
Main example & Compatibility implementation \tabularnewline
\midrule
\code{net10.0}, current package family & \code{net8.0}, older packages from appendix A \tabularnewline
\code{ServeMode.Embedded} & \code{GraphQLToolServeMode.Embedded} inside \code{GraphQLServerOptions} \tabularnewline
\code{QueryContext<T>} with IQueryable \code{With} & Omit that projection; fetch complete rows \tabularnewline
Generated type-class federation attributes & Explicit \code{ObjectType<Speaker>} descriptor configuration \tabularnewline
Injected node-id serializer & Schema service's \code{INodeIdSerializerAccessor}, then its \code{Serializer} property \tabularnewline
\code{TypeSystemConfiguration} & \code{DefinitionBase} \tabularnewline
\code{ObjectTypeConfiguration} & \code{ObjectTypeDefinition} \tabularnewline
\code{DirectiveConfiguration} & \code{DirectiveDefinition} \tabularnewline
\code{Descriptors.Configurations} namespace & \code{Descriptors.Definitions} namespace \tabularnewline
Semantic-non-null export option & Not available in this version \tabularnewline
\bottomrule
\end{tabular}
\end{center}

Batching, cursor connections, global node identifiers and mutation error
payloads remain in this example. The lack of the IQueryable projection
overload does not remove the DataLoader. Parameter names and the
\code{json\_each} form in SQL come from the older EF Core provider, not from
federation.

The generated \code{SessionType} retains its node resolver but is not an
entity in this compatibility schema. \code{SpeakerType} uses descriptor
calls to attach both the key and reference resolver. Keeping a second
generated declaration of that same speaker type would create a collision;
use the complete file set below.

\section{Run the Compatibility Service}

Create a separate example directory and save the files under their shown
paths. Run these commands from \code{src/Sessions}:

\begin{minted}{text}
dotnet restore
dotnet build
dotnet run --no-build -- schema export --output schema.graphql
dotnet run --no-build
\end{minted}

The SDK must be able to build the older target and the corresponding ASP.NET
Core runtime must be installed. Check the emitted schema, not only the build:
\code{Speaker} must carry a resolvable key and appear in \code{\_Entity}.
A clean build with an ignored federation attribute is the failure described
in chapter~\ref{ch:first-subgraph}.

For a local composition, create \code{graph/graph.yaml} relative to that
separate example root:

\begin{minted}{yaml}
version: 1
subgraphs:
  - name: sessions
    routing_url: http://localhost:5001/graphql
    schema:
      file: ../src/Sessions/schema.graphql
\end{minted}

Compose from \code{graph}:

\begin{minted}{text}
wgc router compose -i graph.yaml -o router.json
\end{minted}

The one-service router configuration in chapter~\ref{ch:enter-the-router} can serve this
graph; do not reuse the final four-service graph input from appendix E.

On the direct service port, this entity request reaches the batched speaker
lookup:

\begin{minted}{graphql}
{
  _entities(representations: [
    {__typename: "Speaker", id: "U3BlYWtlcjox"}
  ]) {
    ... on Speaker { id name }
  }
}
\end{minted}

The returned speaker is Ada Fischer, with the same global id. Sending a
Session representation is not a test of this implementation's speaker
resolver: Session is not in this compatibility schema's entity union.

\section{Sessions}

\subsection*{\code{src/Sessions/ConferenceContext.cs}}
\begin{minted}{csharp}
using Microsoft.EntityFrameworkCore;

namespace Sessions;

/// <summary>
/// The conference program, in SQLite. Session carries SpeakerId as a
/// plain column and no navigation property, so EF Core never joins the
/// two tables on its own and every speaker lookup is a resolver's job.
/// </summary>
public sealed class ConferenceContext(
    DbContextOptions<ConferenceContext> options) : DbContext(options)
{
    public DbSet<Session> Sessions => Set<Session>();

    public DbSet<Speaker> Speakers => Set<Speaker>();

    protected override void OnModelCreating(ModelBuilder model)
    {
        // SQLite stores no date type of its own and refuses to sort a
        // DateTimeOffset. Keep the offset out of the column and put it
        // back on the way in; every session here is in conference time.
        model.Entity<Session>()
            .Property(s => s.StartsAt)
            .HasConversion(
                v => v.UtcDateTime,
                v => new DateTimeOffset(v, TimeSpan.Zero));

        model.Entity<Speaker>().HasData(ConferenceData.Speakers);
        model.Entity<Session>().HasData(ConferenceData.Sessions);
    }
}
\end{minted}

\subsection*{\code{src/Sessions/ConferenceData.cs}}
\begin{minted}{csharp}
namespace Sessions;

/// <summary>
/// The conference program, held in memory. The next chapter puts it in
/// SQLite with EF Core; nothing outside this file knows where the data
/// lives.
/// </summary>
public static class ConferenceData
{
    public static readonly IReadOnlyList<Speaker> Speakers =
    [
        new Speaker(1, "Ada Fischer", "Works on query planners."),
        new Speaker(2, "Bruno Kaminski", null),
        new Speaker(3, "Chidi Okafor", "Runs platform engineering.")
    ];

    public static readonly IReadOnlyList<Session> Sessions =
    [
        new Session(
            1,
            "Schemas That Outlive Their Authors",
            "A five-year-old schema, and the decisions that shaped it.",
            1,
            new DateTimeOffset(2026, 9, 14, 9, 0, 0, TimeSpan.Zero),
            45),
        new Session(
            2,
            "Reading a Query Plan",
            null,
            1,
            new DateTimeOffset(2026, 9, 14, 10, 0, 0, TimeSpan.Zero),
            45),
        new Session(
            3,
            "Paging Without Offsets",
            "Why the page you get is not the page you asked for.",
            2,
            new DateTimeOffset(2026, 9, 14, 11, 0, 0, TimeSpan.Zero),
            30),
        new Session(
            4,
            "The Bank That Split Its Graph",
            "Eighteen months, four teams, and the unwritten parts.",
            3,
            new DateTimeOffset(2026, 9, 14, 14, 0, 0, TimeSpan.Zero),
            60)
    ];
}
\end{minted}

\subsection*{\code{src/Sessions/Mutation.cs}}
\begin{minted}{csharp}
using Microsoft.EntityFrameworkCore;

namespace Sessions;

[MutationType]
public static class Mutation
{
    /// <summary>
    /// Moves one session to a new slot, keeping its length.
    /// </summary>
    [Error<SessionNotFoundException>]
    [Error<SpeakerDoubleBookedException>]
    public static async Task<Session> RescheduleSessionAsync(
        [ID<Session>] int sessionId,
        DateTimeOffset startsAt,
        ConferenceContext db,
        CancellationToken ct)
    {
        var session = await db.Sessions
            .FirstOrDefaultAsync(s => s.Id == sessionId, ct)
            ?? throw new SessionNotFoundException(sessionId);

        var ends = startsAt.AddMinutes(session.DurationMinutes);

        // The overlap test is done in memory on purpose. StartsAt goes
        // through a value converter, so SQLite cannot do arithmetic on
        // that column, and one speaker's program is a handful of rows.
        var alsoTheirs = await db.Sessions
            .Where(s => s.SpeakerId == session.SpeakerId)
            .Where(s => s.Id != session.Id)
            .ToListAsync(ct);

        var clash = alsoTheirs.FirstOrDefault(
            s => s.StartsAt < ends
                && startsAt < s.StartsAt.AddMinutes(s.DurationMinutes));

        if (clash is not null)
        {
            throw new SpeakerDoubleBookedException(
                session.SpeakerId,
                clash.Id);
        }

        var moved = session with { StartsAt = startsAt };
        db.Entry(session).CurrentValues.SetValues(moved);
        await db.SaveChangesAsync(ct);

        return moved;
    }
}
\end{minted}

\subsection*{\code{src/Sessions/Program.cs}}
\begin{minted}{csharp}
using HotChocolate.AspNetCore;
using Microsoft.EntityFrameworkCore;
using Sessions;

var builder = WebApplication.CreateBuilder(args);

// EF Core reports every command through the logging pipeline as well,
// at Information and with a duration attached. StatementLog below is
// this service's instrument, so silence the other one rather than have
// every statement printed twice.
builder.Logging.AddFilter(
    "Microsoft.EntityFrameworkCore.Database.Command",
    LogLevel.Warning);

builder.Services.AddDbContextFactory<ConferenceContext>(options =>
{
    options.UseSqlite("Data Source=conference.db");

    if (builder.Environment.IsDevelopment())
    {
        options.AddInterceptors(new StatementLog());
    }
});

builder
    .AddGraphQL()
    .AddSessionsTypes()
    .RegisterDbContextFactory<ConferenceContext>()
    .AddQueryContext()
    .AddPagingArguments()
    .AddGlobalObjectIdentification()
    .AddMutationConventions()
    .AddApolloFederation()
    .ModifyCostOptions(options => options.ApplyCostDefaults = false)
    .TryAddTypeInterceptor<ShareNodeFields>();

var app = builder.Build();

using (var scope = app.Services.CreateScope())
{
    var factory = scope.ServiceProvider
        .GetRequiredService<IDbContextFactory<ConferenceContext>>();
    using var db = factory.CreateDbContext();
    db.Database.EnsureCreated();
}

app.MapGraphQL()
    .WithOptions(new GraphQLServerOptions
    {
        Tool = { ServeMode = GraphQLToolServeMode.Embedded }
    });

app.RunWithGraphQLCommands(args);
\end{minted}

\subsection*{\code{src/Sessions/Properties/launchSettings.json}}
\begin{minted}[fontsize=\footnotesize,breakanywhere]{json}
{
  "$schema": "https://json.schemastore.org/launchsettings.json",
  "profiles": {
    "http": {
      "commandName": "Project",
      "dotnetRunMessages": true,
      "launchBrowser": false,
      "applicationUrl": "http://localhost:5001",
      "environmentVariables": {
        "ASPNETCORE_ENVIRONMENT": "Development"
      }
    }
  }
}
\end{minted}

\subsection*{\code{src/Sessions/Query.cs}}
\begin{minted}{csharp}
using Microsoft.EntityFrameworkCore;

namespace Sessions;

[QueryType]
public static class Query
{
    /// <summary>
    /// Sessions on the schedule, earliest first.
    /// </summary>
    [UsePaging(DefaultPageSize = 2, MaxPageSize = 10)]
    public static IQueryable<Session> GetSessions(ConferenceContext db)
        => db.Sessions.OrderBy(s => s.StartsAt);

    /// <summary>
    /// One session by its identifier, or null if no session has it.
    /// </summary>
    [GraphQLDeprecated(
        "Ask for `node(id:)` instead. This field takes the database "
        + "key, which is only unique among sessions.")]
    public static async Task<Session?> GetSessionByIdAsync(
        int id,
        ConferenceContext db,
        CancellationToken ct)
        => await db.Sessions
            .FirstOrDefaultAsync(s => s.Id == id, ct);
}
\end{minted}

\subsection*{\code{src/Sessions/Session.cs}}
\begin{minted}{csharp}
namespace Sessions;

/// <summary>
/// A talk on the conference program.
/// </summary>
/// <param name="Id">Stable identifier for this session.</param>
/// <param name="Title">The line that goes in the program.</param>
/// <param name="Abstract">The longer text, if there is one.</param>
/// <param name="SpeakerId">Who gives it. Not part of the schema.</param>
/// <param name="StartsAt">When it starts, in conference time.</param>
/// <param name="DurationMinutes">How long the slot is.</param>
public sealed record Session(
    int Id,
    string Title,
    string? Abstract,
    [property: GraphQLIgnore] int SpeakerId,
    DateTimeOffset StartsAt,
    int DurationMinutes);
\end{minted}

\subsection*{\code{src/Sessions/SessionNotFoundException.cs}}
\begin{minted}{csharp}
using HotChocolate.Types.Relay;

namespace Sessions;

/// <summary>
/// No session on the program carries the identifier that was asked for.
/// </summary>
public sealed class SessionNotFoundException(int sessionId)
    : Exception($"No session has id {sessionId}.")
{
    /// <summary>
    /// The identifier that matched nothing.
    /// </summary>
    [ID<Session>]
    public int SessionId { get; } = sessionId;
}
\end{minted}

\subsection*{\code{src/Sessions/SessionType.cs}}
\begin{minted}{csharp}
using Microsoft.EntityFrameworkCore;

namespace Sessions;

[ObjectType<Session>]
public static partial class SessionType
{
    /// <summary>
    /// The person giving this session.
    /// </summary>
    public static async Task<Speaker?> GetSpeakerAsync(
        [Parent(nameof(Session.SpeakerId))] Session session,
        ISpeakerByIdDataLoader speakerById,
        CancellationToken ct)
        => await speakerById.LoadAsync(session.SpeakerId, ct);

    /// <summary>
    /// Loads one session by the identifier inside a global node id.
    /// </summary>
    // Chapter 4's callout tells a reader on 14 to drop the
    // QueryContext parameter and the .With(query) line, because 14 has
    // no IQueryable overload of With. This is that file, so it is
    // written the way the callout says rather than the way a resolver
    // starting from scratch here would be.
    [NodeResolver]
    public static async Task<Session?> GetSessionAsync(
        int id,
        ConferenceContext db,
        CancellationToken ct)
        => await db.Sessions
            .Where(s => s.Id == id)
            .FirstOrDefaultAsync(ct);
}
\end{minted}

\subsection*{\code{src/Sessions/Sessions.csproj}}
\begin{minted}[fontsize=\footnotesize,breakanywhere]{xml}
<Project Sdk="Microsoft.NET.Sdk.Web">

  <PropertyGroup>
    <TargetFramework>net8.0</TargetFramework>
    <Nullable>enable</Nullable>
    <ImplicitUsings>enable</ImplicitUsings>
    <GenerateDocumentationFile>true</GenerateDocumentationFile>
    <NoWarn>$(NoWarn);CS1591</NoWarn>
  </PropertyGroup>

  <ItemGroup>
    <PackageReference Include="HotChocolate.ApolloFederation" Version="14.3.1" />
    <PackageReference Include="HotChocolate.AspNetCore" Version="14.3.1" />
    <PackageReference Include="HotChocolate.AspNetCore.CommandLine" Version="14.3.1" />
    <PackageReference Include="HotChocolate.Data.EntityFramework" Version="14.3.1" />
    <PackageReference Include="HotChocolate.Types.Analyzers" Version="14.3.1" />
    <PackageReference Include="Microsoft.EntityFrameworkCore.Sqlite" Version="9.0.19" />
  </ItemGroup>

</Project>
\end{minted}

\subsection*{\code{src/Sessions/ShareNodeFields.cs}}
\begin{minted}{csharp}
using HotChocolate.Configuration;
using HotChocolate.Language;
using HotChocolate.Types.Descriptors.Definitions;

namespace Sessions;

/// <summary>
/// Applies @shareable to the two fields global object identification
/// generates, and to nothing else. The same class as on 16, with the
/// four type names 16 renamed: DefinitionBase for
/// TypeSystemConfiguration, ObjectTypeDefinition for
/// ObjectTypeConfiguration, DirectiveDefinition for
/// DirectiveConfiguration, and the Definitions namespace for
/// Configurations.
/// </summary>
internal sealed class ShareNodeFields : TypeInterceptor
{
    public override void OnBeforeCompleteType(
        ITypeCompletionContext context,
        DefinitionBase definition)
    {
        if (definition is not ObjectTypeDefinition type
            || type.Name != "Query")
        {
            return;
        }

        foreach (var field in type.Fields)
        {
            if (field.Name is "node" or "nodes")
            {
                field.Directives.Add(
                    new DirectiveDefinition(
                        new DirectiveNode("shareable")));
            }
        }
    }
}
\end{minted}

\subsection*{\code{src/Sessions/Speaker.cs}}
\begin{minted}{csharp}
namespace Sessions;

/// <summary>
/// A person giving one or more sessions.
/// </summary>
/// <param name="Id">Stable identifier for this speaker.</param>
/// <param name="Name">The name to print in the program.</param>
/// <param name="Bio">A sentence or two, if the speaker sent any.</param>
public sealed record Speaker(int Id, string Name, string? Bio);
\end{minted}

\subsection*{\code{src/Sessions/SpeakerDataLoader.cs}}
\begin{minted}{csharp}
using Microsoft.EntityFrameworkCore;

namespace Sessions;

public static class SpeakerDataLoader
{
    /// <summary>
    /// Every speaker a batch of sessions asked for, keyed by identifier.
    /// </summary>
    [DataLoader]
    public static async Task<Dictionary<int, Speaker>>
        GetSpeakerByIdAsync(
            IReadOnlyList<int> ids,
            ConferenceContext db,
            CancellationToken ct)
        => await db.Speakers
            .Where(s => ids.Contains(s.Id))
            .ToDictionaryAsync(s => s.Id, ct);
}
\end{minted}

\subsection*{\code{src/Sessions/SpeakerDoubleBookedException.cs}}
\begin{minted}{csharp}
using HotChocolate.Types.Relay;

namespace Sessions;

/// <summary>
/// The new slot overlaps another session the same speaker gives.
/// </summary>
public sealed class SpeakerDoubleBookedException(
    int speakerId,
    int clashingSessionId)
    : Exception(
        $"Speaker {speakerId} already gives session "
        + $"{clashingSessionId} in that slot.")
{
    /// <summary>
    /// The speaker who cannot be in two rooms at once.
    /// </summary>
    [ID<Speaker>]
    public int SpeakerId { get; } = speakerId;

    /// <summary>
    /// The session already occupying the slot.
    /// </summary>
    [ID<Session>]
    public int ClashingSessionId { get; } = clashingSessionId;
}
\end{minted}

\subsection*{\code{src/Sessions/SpeakerType.cs}}
\begin{minted}{csharp}
using HotChocolate.ApolloFederation.Types;
using HotChocolate.Resolvers;
using HotChocolate.Types.Relay;
using Microsoft.Extensions.DependencyInjection;

namespace Sessions;

/// <summary>
/// Speaker as a federated entity on HotChocolate 14. The attribute
/// idiom the generated type-class example uses does not work here:
/// 14's source generator has no mechanism for applying an arbitrary
/// descriptor attribute declared on the type class, so [Key] and
/// [ReferenceResolver] compile and are never read. Configuring the type
/// the code-first way applies both, because the descriptor calls are
/// made rather than discovered.
///
/// The cost is that this type leaves the generator entirely: there is
/// no [ObjectType&lt;Speaker&gt;] partial class for Speaker in this
/// version, because two declarations of one type collide.
/// </summary>
public sealed class SpeakerType : ObjectType<Speaker>
{
    protected override void Configure(
        IObjectTypeDescriptor<Speaker> descriptor)
    {
        descriptor
            .ImplementsNode()
            .ResolveNodeWith(
                typeof(SpeakerType).GetMethod(
                    nameof(ResolveByKeyAsync))!);

        descriptor
            .Key("id")
            .ResolveReferenceWith(
                _ => ResolveReferenceAsync(default!, default!));
    }

    public static async Task<Speaker?> ResolveByKeyAsync(
        IResolverContext context,
        int id)
    {
        var loader = context.Service<ISpeakerByIdDataLoader>();
        return await loader.LoadAsync(id, context.RequestAborted);
    }

    /// <summary>
    /// The reference resolver. The key arrives as the undecoded node id
    /// string here exactly as it does on 16, and the decode is the same
    /// call; only the way the method is attached to the type differs.
    /// </summary>
    public static async Task<Speaker?> ResolveReferenceAsync(
        string id,
        IResolverContext context)
    {
        var serializer = context.Schema.Services
            .GetRequiredService<INodeIdSerializerAccessor>()
            .Serializer;
        var nodeId = serializer.Parse(id, typeof(int));

        if (nodeId.TypeName != nameof(Speaker))
        {
            return null;
        }

        var loader = context.Service<ISpeakerByIdDataLoader>();
        return await loader.LoadAsync(
            (int)nodeId.InternalId!, context.RequestAborted);
    }
}
\end{minted}

\subsection*{\code{src/Sessions/StatementLog.cs}}
\begin{minted}{csharp}
using System.Data.Common;
using Microsoft.EntityFrameworkCore.Diagnostics;

namespace Sessions;

/// <summary>
/// Prints every statement the service sends to SQLite, numbered, so a
/// request's cost can be read off the console. The counter is
/// process-wide and never reset: start the service, send one request,
/// and the last number is what that request cost.
/// </summary>
public sealed class StatementLog : DbCommandInterceptor
{
    private static int _count;

    public override InterceptionResult<DbDataReader> ReaderExecuting(
        DbCommand command,
        CommandEventData eventData,
        InterceptionResult<DbDataReader> result)
    {
        Print(command);
        return result;
    }

    public override ValueTask<InterceptionResult<DbDataReader>>
        ReaderExecutingAsync(
            DbCommand command,
            CommandEventData eventData,
            InterceptionResult<DbDataReader> result,
            CancellationToken cancellationToken = default)
    {
        Print(command);
        return ValueTask.FromResult(result);
    }

    private static void Print(DbCommand command)
    {
        var n = Interlocked.Increment(ref _count);
        Console.WriteLine($"-- statement {n}");
        Console.WriteLine(command.CommandText);
    }
}
\end{minted}

\subsection*{\code{src/Sessions/appsettings.Development.json}}
\begin{minted}[fontsize=\footnotesize,breakanywhere]{json}
{
  "Logging": {
    "LogLevel": {
      "Default": "Information",
      "Microsoft.AspNetCore": "Warning"
    }
  }
}
\end{minted}

\subsection*{\code{src/Sessions/appsettings.json}}
\begin{minted}[fontsize=\footnotesize,breakanywhere]{json}
{
  "Logging": {
    "LogLevel": {
      "Default": "Information",
      "Microsoft.AspNetCore": "Warning"
    }
  },
  "AllowedHosts": "*"
}
\end{minted}
